\item \model{Nemotron 3}

\paragraph*{حذف مطلق \en{RoPE} در لایه‌های توجه با تکیه بر بازنمایی ضمنی \en{Mamba-2}}
در معماری مدل‌های خانواده \model{NVIDIA Nemotron 3} \cite{nvidia2025nemotron3} (در نسخه‌های \en{Nano}، \en{Super} و \en{Ultra})، پیشرفتی بنیادین در زمینه کدگذاری موقعیت رخ داده است: \textbf{تمام لایه‌های توجه بدون کدگذاری دورانی موقعیت (\en{No RoPE}) طراحی شده‌اند}.
به طور سنتی، برون‌یابی فرکانس‌های زاویه‌ای RoPE فراتر از طول آموزش اولیه، با افت شدید توجه و واگرایی ضرب داخلی توکن‌ها همراه است. در Nemotron 3، ترانسفورمر استاندارد با لایه‌های مدل فضای حالت \en{Mamba-2} ترکیب شده است. لایه‌های بازگشتی Mamba-2 ساختار محاسباتی پیوسته به صورت زیر دارند:
\begin{equation}
    h_t = \mathbf{A}_t h_{t-1} + \mathbf{B}_t x_t, \quad y_t = \mathbf{C}_t h_t
\end{equation}
از آنجا که ماتریس انتقال $\mathbf{A}_t$ و ترتیب عملیات کانولوشن موضعی ذاتاً توالی زمانی، علیت و جایگاه گام‌های ورودی را در بردار حالت پنهان $h_t$ ثبت می‌کنند، اطلاعات زمانی پیش از ورود به لایه‌های توجه، در تانسور ورودی تعبیه شده است.

\paragraph*{پیامدهای تئوریک بر اکستنشن زمینه تا مرز ۱ میلیون توکن}
عدم استفاده از RoPE مزایای چندگانه‌ای به همراه داشته است:
\begin{itemize}
    \item \textbf{ریشه‌کنی خطای توزیع خارج از دامنه (\en{OOD RoPE Issues}):} مدل در مرحله گسترش بافت نیازی به دستکاری فرکانس مبنا یا درون‌یابی‌های زاویه‌ای ندارد.
    \item \textbf{انطباق با محاسبات ممیز شناور ۴ بیتی (\en{NVFP4}):} خطاهای گردکردن ناشی از ضرب‌های ماتریس دوران در فرمت‌های کم‌بیت به طور کامل حذف می‌گردد.
    \item \textbf{تداوم شیب درست‌نمایی منفی (\en{NLL}):} منحنی‌های ارزیابی تجربی بر روی کدهای منبع با طول بیش از یک میلیون توکن نشان داد که منفی لگاریتم درست‌نمایی تجمعی به‌صورت یکنواخت با افزایش طول توالی از ۱۲۸ تا ۱M توکن بر اساس تابع توانی ($R^2 = 0.883$) بهبود می‌یابد (شکل ۶ مقاله).
\end{itemize}